\section{Conclusions}
\label{sec:conclusion}

Public laboratory-automation discussions contain enough technical depth to support a rigorous qualitative pilot and a bounded quantitative program. Their strongest scientific use is measurement design. They expose which units can be defined, which fields are missing, where evidence fails to support a claim, and which operational outcomes a prospective study should record.

The study identifies ten bottlenecks organized across ecosystem knowledge, interfaces and representations, runtime coordination, evaluation, and AI context. The deeper structure concerns boundaries. Knowledge must become a maintained artifact. Device access must become an executable and testable interface. Source must become an identified deployment. Resource data must become a validated physical model. Observations must become linked evidence. Interrupted execution must become a verified recovery decision. Generated methods must become experimentally evaluated procedures.

The strongest next empirical study concerns observability and recovery. The current pilot finds the strongest association between these constructs and supplies direct incident mechanisms. A prospective incident dataset can test whether command--state--material linkage predicts time to a verified safe state, material salvage, final disposition completeness, and recurrence. The public artifacts developed here provide an immediate contribution while enabling that future research. The troubleshooting schema improves new questions. The resolved-knowledge index preserves bounded answers. The interface, resource, event, and scheduling registries create the denominators required for stronger operational science.
